\documentclass{amsart}
% For dealing with references we use the comment environment
\usepackage{verbatim}
\newenvironment{reference}{\comment}{\endcomment}
%\newenvironment{reference}{}{}
\newenvironment{slogan}{\comment}{\endcomment}
\newenvironment{history}{\comment}{\endcomment}

% For commutative diagrams we use Xy-pic
\usepackage[all]{xy}

% We use 2cell for 2-commutative diagrams.
\xyoption{2cell}
\UseAllTwocells

% We use multicol for the list of chapters between chapters
\usepackage{multicol}

% This is generally recommended for better output
\usepackage{lmodern}
\usepackage[T1]{fontenc}

% For cross-file-references

% Package for hypertext links:
\usepackage{hyperref}

% For any local file, say "hello.tex" you want to link to please

% Theorem environments.
%
\theoremstyle{plain}
\newtheorem{theorem}[subsection]{Theorem}
\newtheorem{proposition}[subsection]{Proposition}
\newtheorem{lemma}[subsection]{Lemma}

\theoremstyle{definition}
\newtheorem{definition}[subsection]{Definition}
\newtheorem{example}[subsection]{Example}
\newtheorem{exercise}[subsection]{Exercise}
\newtheorem{situation}[subsection]{Situation}

\theoremstyle{remark}
\newtheorem{remark}[subsection]{Remark}
\newtheorem{remarks}[subsection]{Remarks}

\numberwithin{equation}{subsection}

% Macros
%
\def\lim{\mathop{\mathrm{lim}}\nolimits}
\def\colim{\mathop{\mathrm{colim}}\nolimits}
\def\Spec{\mathop{\mathrm{Spec}}}
\def\Hom{\mathop{\mathrm{Hom}}\nolimits}
\def\Ext{\mathop{\mathrm{Ext}}\nolimits}
\def\SheafHom{\mathop{\mathcal{H}\!\mathit{om}}\nolimits}
\def\SheafExt{\mathop{\mathcal{E}\!\mathit{xt}}\nolimits}
\def\Sch{\mathit{Sch}}
\def\Mor{\mathop{\mathrm{Mor}}\nolimits}
\def\Ob{\mathop{\mathrm{Ob}}\nolimits}
\def\Sh{\mathop{\mathit{Sh}}\nolimits}
\def\NL{\mathop{N\!L}\nolimits}
\def\CH{\mathop{\mathrm{CH}}\nolimits}
\def\proetale{{pro\text{-}\acute{e}tale}}
\def\etale{{\acute{e}tale}}
\def\QCoh{\mathit{QCoh}}
\def\Ker{\mathop{\mathrm{Ker}}}
\def\Im{\mathop{\mathrm{Im}}}
\def\Coker{\mathop{\mathrm{Coker}}}
\def\Coim{\mathop{\mathrm{Coim}}}

% Boxtimes
%
\DeclareMathSymbol{\boxtimes}{\mathbin}{AMSa}{"02}

%
% Macros for moduli stacks/spaces
%
\def\QCohstack{\mathcal{QC}\!\mathit{oh}}
\def\Cohstack{\mathcal{C}\!\mathit{oh}}
\def\Spacesstack{\mathcal{S}\!\mathit{paces}}
\def\Quotfunctor{\mathrm{Quot}}
\def\Hilbfunctor{\mathrm{Hilb}}
\def\Curvesstack{\mathcal{C}\!\mathit{urves}}
\def\Polarizedstack{\mathcal{P}\!\mathit{olarized}}
\def\Complexesstack{\mathcal{C}\!\mathit{omplexes}}
% \Pic is the operator that assigns to X its picard group, usage \Pic(X)
% \Picardstack_{X/B} denotes the Picard stack of X over B
% \Picardfunctor_{X/B} denotes the Picard functor of X over B
\def\Pic{\mathop{\mathrm{Pic}}\nolimits}
\def\Picardstack{\mathcal{P}\!\mathit{ic}}
\def\Picardfunctor{\mathrm{Pic}}
\def\Deformationcategory{\mathcal{D}\!\mathit{ef}}

\setlength{\emergencystretch}{3em}
\begin{document}
\title[Stacks Project, Tag 0F0V]{288 --- Affine etale cohomological dimension is bounded by dimension plus the field cohomological dimension}
\maketitle
\noindent Copyright (C) 2005--2025 Johan de Jong. Stacks Project authors. Source: \href{https://stacks.math.columbia.edu/tag/0F0V}{Tag 0F0V}.

\noindent Editorial difficulty note (not author text). Difficulty H3: Noether normalization alone does not give the desired bound: the proof needs local strict-henselian estimates indexed by residue transcendence degree. Two support estimates control the boundary of a partial compactification, and a careful two-stage Leray analysis bounds every diagonal of the spectral sequence..

\noindent Editorial provenance note (not author text). History: excerpt from revision \texttt{a0444\allowbreak 6e57e\allowbreak c1fbc\allowbreak 252a8\allowbreak 71afc\allowbreak ec775\allowbreak 2fb28\allowbreak 07b14}, etale-cohomology.tex, lines 19762--19913. Reference presentation was adapted; original-author context and core support blocks were retained or added. The mathematical wording is unchanged. Licensed under GNU FDL 1.2 or later; no invariant sections or cover texts. See ../licenses/COPYING. Context blocks reproduce author definitions and supporting results; they are not additional counted problems.

\begin{proposition}
\label{proposition-cd-affine}
Let $K$ be a field. Let $X$ be an affine scheme of finite type over $K$.
Then we have $\text{cd}(X) \leq \dim(X) + \text{cd}(K)$.
\end{proposition}

\begin{proof}
We will prove this by induction on $\dim(X)$.
Let $\mathcal{F}$ be an abelian torsion sheaf on $X_\etale$.

\medskip\noindent
The case $\dim(X) = 0$. In this case the structure morphism
$f : X \to \Spec(K)$ is finite. Hence we see that $R^if_*\mathcal{F} = 0$
for $i > 0$, see Proposition \href{https://stacks.math.columbia.edu/tag/03QP}{proposition finite higher direct image zero (Tag 03QP)}.
Thus $H^i_\etale(X, \mathcal{F}) = H^i_\etale(\Spec(K), f_*\mathcal{F})$
by the Leray spectral sequence for $f$
(Cohomology on Sites, Lemma \href{https://stacks.math.columbia.edu/tag/0732}{lemma Leray (Tag 0732)})
and the result is clear.

\medskip\noindent
The case $\dim(X) = 1$. This is Lemma \href{https://stacks.math.columbia.edu/tag/0F0S}{lemma cd curve over field (Tag 0F0S)}.

\medskip\noindent
Assume $d = \dim(X) > 1$ and the proposition holds for finite type
affine schemes of dimension $< d$ over fields.
By Noether normalization, see for example
Varieties, Lemma \href{https://stacks.math.columbia.edu/tag/0CBI}{lemma noether normalization affine (Tag 0CBI)},
there exists a finite morphism $f : X \to \mathbf{A}^d_K$.
Recall that $R^if_*\mathcal{F} = 0$ for $i > 0$ by
Proposition \href{https://stacks.math.columbia.edu/tag/03QP}{proposition finite higher direct image zero (Tag 03QP)}.
By the Leray spectral sequence for $f$
(Cohomology on Sites, Lemma \href{https://stacks.math.columbia.edu/tag/0732}{lemma Leray (Tag 0732)})
we conclude that it suffices to prove the result for $\pi_*\mathcal{F}$
on $\mathbf{A}^d_K$.

\medskip\noindent
Interlude I. Let $j : X \to Y$ be an open immersion of smooth
$d$-dimensional varieties over $K$ (not necessarily affine)
whose complement is the support of an effective Cartier divisor $D$.
The sheaves $R^qj_*\mathcal{F}$ for $q > 0$ are supported on $D$.
We claim that $(R^qj_*\mathcal{F})_{\overline{y}} = 0$
for $a = \text{trdeg}_K(\kappa(y)) > d - q$. Namely, by
Theorem \href{https://stacks.math.columbia.edu/tag/03Q9}{theorem higher direct images (Tag 03Q9)} we have
$$
(R^qj_*\mathcal{F})_{\overline{y}} =
H^q(\Spec(\mathcal{O}_{Y, y}^{sh}) \times_Y X, \mathcal{F})
$$
Choose a local equation $f \in \mathfrak m_y \subset \mathcal{O}_{Y, y}$
for $D$. Then we have
$$
\Spec(\mathcal{O}_{Y, y}^{sh}) \times_Y X =
\Spec(\mathcal{O}_{Y, y}^{sh}[1/f])
$$
Using Lemma \href{https://stacks.math.columbia.edu/tag/0F0U}{lemma strictly henselian (Tag 0F0U)} we get an embedding
$$
K(t_1, \ldots, t_a)^{sep}(x) =
K(t_1, \ldots, t_a)^{sep}[x]_{(x)}[1/x]
\longrightarrow
\mathcal{O}_{Y, y}^{sh}[1/f]
$$
Since the transcendence degree over $K$ of the fraction field of
$\mathcal{O}_{Y, y}^{sh}$ is $d$, we see that $\mathcal{O}_{Y, y}^{sh}[1/f]$
is a filtered colimit of $(d - a - 1)$-dimensional finite type algebras over
the field $K(t_1, \ldots, t_a)^{sep}(x)$ which itself has cohomological
dimension $1$ by Lemma \href{https://stacks.math.columbia.edu/tag/0F0T}{lemma cd field extension (Tag 0F0T)}. Thus by induction
hypothesis and Lemma \href{https://stacks.math.columbia.edu/tag/0F0R}{lemma cd limit (Tag 0F0R)}
we obtain the desired vanishing.

\medskip\noindent
Interlude II. Let $Z$ be a smooth variety over $K$ of dimension $d - 1$.
Let $E_a \subset Z$ be the set of points $z \in Z$ with
$\text{trdeg}_K(\kappa(z)) \leq a$. Observe that $E_a$ is closed
under specialization, see
Varieties, Lemma \href{https://stacks.math.columbia.edu/tag/0A21}{lemma dimension locally algebraic (Tag 0A21)}.
Suppose that $\mathcal{G}$ is a torsion abelian sheaf on $Z$
whose support is contained in $E_a$. Then we claim that
$H^b_\etale(Z, \mathcal{G}) = 0$ for $b > a + \text{cd}(K)$.
Namely, we can write $\mathcal{G} = \colim \mathcal{G}_i$
with $\mathcal{G}_i$ a torsion abelian sheaf
supported on a closed subscheme $Z_i$ contained in $E_a$, see
Lemma \href{https://stacks.math.columbia.edu/tag/0F0N}{lemma support in subset (Tag 0F0N)}.
Then the induction hypothesis kicks in to imply
the desired vanishing for $\mathcal{G}_i$\footnote{Here
we first use Proposition \href{https://stacks.math.columbia.edu/tag/04CA}{proposition closed immersion pushforward (Tag 04CA)}
to write $\mathcal{G}_i$ as the pushforward of a sheaf on $Z_i$,
the induction hypothesis gives the vanishing for this sheaf on $Z_i$, and
the Leray spectral sequence for $Z_i \to Z$ gives the vanishing
for $\mathcal{G}_i$.}. Finally, we
conclude by Theorem \href{https://stacks.math.columbia.edu/tag/09YQ}{theorem colimit (Tag 09YQ)}.

\medskip\noindent
Consider the commutative diagram
$$
\xymatrix{
\mathbf{A}^d_K \ar[rd]_f \ar[rr]_-j & &
\mathbf{P}^1_K \times_K \mathbf{A}^{d - 1}_K \ar[ld]^g \\
& \mathbf{A}^{d - 1}_K
}
$$
Observe that $j$ is an open immersion of smooth $d$-dimensional
varieties whose complement is an effective Cartier divisor $D$. Thus
we may use the results obtained in interlude I.
We are going to study the relative Leray spectral sequence
$$
E_2^{p, q} = R^pg_*R^qj_*\mathcal{F} \Rightarrow R^{p + q}f_*\mathcal{F}
$$
Since $R^qj_*\mathcal{F}$ for $q > 0$ is supported on $D$ and since
$g|_D : D \to \mathbf{A}^{d - 1}_K$ is an isomorphism, we find
$R^pg_*R^qj_*\mathcal{F} = 0$ for $p > 0$ and $q > 0$. Moreover, we have
$R^qj_*\mathcal{F} = 0$ for $q > d$. On the other hand, $g$ is a
proper morphism of relative dimension $1$. Hence by
Lemma \href{https://stacks.math.columbia.edu/tag/095U}{lemma cohomological dimension proper (Tag 095U)}
we see that $R^pg_*j_*\mathcal{F} = 0$ for $p > 2$. Thus the $E_2$-page
of the spectral sequence looks like this
$$
\begin{matrix}
g_*R^dj_*\mathcal{F} & 0 & 0 \\
\ldots & \ldots & \ldots \\
g_*R^2j_*\mathcal{F} & 0 & 0 \\
g_*R^1j_*\mathcal{F} & 0 & 0 \\
g_*j_*\mathcal{F} & R^1g_*j_*\mathcal{F} & R^2g_*j_*\mathcal{F}
\end{matrix}
$$
We conclude that $R^qf_*\mathcal{F} = g_*R^qj_*\mathcal{F}$ for $q > 2$.
By interlude I we see that the support of $R^qf_*\mathcal{F}$ for $q > 2$
is contained in the set of points of $\mathbf{A}^{d - 1}_K$
whose residue field has transcendence degree $\leq d - q$.
By interlude II
$$
H^p(\mathbf{A}^{d - 1}_K, R^qf_*\mathcal{F}) = 0
\text{ for }p > d - q + \text{cd}(K)\text{ and }q > 2
$$
On the other hand, by Theorem \href{https://stacks.math.columbia.edu/tag/03Q9}{theorem higher direct images (Tag 03Q9)}
we have $R^2f_*\mathcal{F}_{\overline{\eta}} =
H^2(\mathbf{A}^1_{\overline{\eta}}, \mathcal{F}) = 0$
(vanishing by the case of dimension $1$) where $\eta$ is the
generic point of $\mathbf{A}^{d - 1}_K$.
Hence by interlude II again we see
$$
H^p(\mathbf{A}^{d - 1}_K, R^2f_*\mathcal{F}) = 0
\text{ for }p > d - 2 + \text{cd}(K)
$$
Finally, we have
$$
H^p(\mathbf{A}^{d - 1}_K, R^qf_*\mathcal{F}) = 0
\text{ for }p > d - 1 + \text{cd}(K)\text{ and }q = 0, 1
$$
by induction hypothesis. Combining everything we just said
with the Leray spectral sequence
$H^p(\mathbf{A}^{d - 1}_K, R^qf_*\mathcal{F}) \Rightarrow
H^{p + q}(\mathbf{A}^d_K, \mathcal{F})$ we conclude.
\end{proof}
\end{document}
